\section{Finite boundaries for the chapter 5.4 comparison}

\begin{annexitem}{annex:p19.r54.disabler-boundary}{thm:p19-r54-disabler-boundary}{The unchanged material contradiction and the explicitly altered rule}
\begin{proof}
Any valuation satisfying p and \(p\to q\) satisfies q,
so cannot also satisfy \(\neg q\). The initial base has
no model. Any superset of a model-free premise set is still
model-free, since a model of the superset would satisfy its
subset. This proves that adding r, \(r\to\neg q\), or
any other formulas does not restore its consistency.
For the changed rule take \(p=1,q=0,r=1\). Then
\(p\land\neg r\) is false and its implication to q true;
p, \(\neg q\) and r are all true. The same valuation also
satisfies \(r\to\neg q\). Thus the altered base is
consistent, with or without this extra disabling rule.
It does not satisfy the original \(p\to q\), so the
alteration is not an equivalence-preserving rewriting of
that earlier premise. The reasoning can model a default
interpretation with an exception, but it supplies no observation
that a participant actually used this representation.
\end{proof}
\end{annexitem}

\begin{annexitem}{annex:p19.r54.probability-boundary}{thm:p19-r54-probability-boundary}{The four weighted possibilities and every denominator condition}
\begin{proof}
Normalize the four stipulated weights by their total 10.
The probabilities of \(p\land q,p\land\neg q,
\neg p\land q,\neg p\land\neg q\) are then
\(3/10,1/10,2/10,4/10\). Consequently
\(\Pr(p)=4/10\), \(\Pr(q)=5/10\), and
\(\Pr(p\land q)=3/10\), yielding
\(\Pr(q\mid p)=3/4\) and \(\Pr(p\mid q)=3/5\).
Material \(p\to q\) fails only in the second cell,
so has probability \(1-1/10=9/10\), not 3/4.
This calculation distinguishes conditional confidence from
the probability of a particular truth-functional proposition.

For any distribution, \(\Pr(\neg q)=1\) implies
\(\Pr(q)=0\). Since \(p\land q\) is a subset of q,
\(\Pr(p\land q)=0\). Dividing by \(\Pr(p)>0\)
gives \(\Pr(q\mid p)=0\). If \(\Pr(p)=0\), the
ordinary quotient is \(0/0\), which defines no value.
An extended conditional-probability convention would require
separate specification. If instead \(\Pr(p)=1\),
\(\Pr(\neg p)=0\), and splitting q into its p and
not-p parts gives \(\Pr(p\land q)=\Pr(q)\).
Dividing by 1 gives the stated probability equality.

Ordinary conditionalization on an event E defines
\(\Pr_E(B)=\Pr(B\land E)/\Pr(E)\), so needs
\(\Pr(E)>0\). In the illustrative distribution,
conditioning on \(\neg q\) yields probabilities
\(1/5,4/5\) for its two possible cells; thus the new
\(\Pr_{\neg q}(p)=1/5>0\) and its q-given-p is zero.
Conditioning on p gives the two probabilities 3/4 and 1/4,
so its new \(\Pr_p(q)=\Pr_p(q\mid p)=3/4\).
These are operations on stipulated distributions, not measured
changes in participants' credences. If the initial state instead
assigned q probability 1, learning \(\neg q\) would not be
handled by this ordinary zero-denominator operation. Nothing
in the proof identifies response proportions with personal probabilities.
\end{proof}
\end{annexitem}

\begin{annexitem}{annex:p19.r54.two-revisions}{thm:p19-r54-two-revisions}{All eight postulates for the two fixed-prior finite operations}
\begin{proof}
Let \(W=\{(0,0),(0,1),(1,0),(1,1)\}\), with entries
giving p and q. Put \(v_0=(1,1)\). The material base
\(p\to q,p\) has exactly this model, so
\(K=\operatorname{Th}(\{v_0\})\).
For either displayed total order and any formula A, let
\(m_A=\min([A])\) and define \(R_A=\operatorname{Th}(m_A)\).
Because W is finite and the order strict, \(m_A\) is a
singleton when \([A]\ne\varnothing\), and is empty otherwise.
Each subset of W is definable by a finite disjunction of
valuation cubes. Hence its full theory has exactly that subset
as models. In particular expanding a theory of X by B has
model set \(X\cap[B]\), also when X is empty.

We check the eight revision postulates for this fixed K,
for arbitrary A,B in the finite language, separately from
any psychological interpretation of the orders.

Closure: any classical consequence of formulas true in every
member of \(m_A\) is true in every such member, so
\(R_A\) is deductively closed. Success: each member of
\(m_A\) satisfies A, so \(A\in R_A\); when the minimum
is empty the inconsistent theory contains A as well.
For Inclusion, if \(v_0\in[A]\), it is the unique minimum,
and both \(K+A\) and \(R_A\) are its theory. If
\(v_0\notin[A]\), \(K+A\) is inconsistent and contains
every formula, so \(R_A\subseteq K+A\). Thus Inclusion
is unconditional. For Vacuity, \(\neg A\notin K\)
means \(v_0\models A\), since K is this valuation's
complete theory. The same equality of expansion and revision
then gives the required reverse inclusion.

Consistency: a consistent formula in this finite language has
a valuation, so \(m_A\) is nonempty and its theory is
consistent. Extensionality: classically equivalent inputs have
the same valuation set and hence the same selected minimum.
For Superexpansion observe
\(m_A\cap[B]\subseteq m_{A\land B}\).
If the left side has a member, it is already the least
A-valuation and also a B-valuation, so no A-and-B valuation
precedes it; if it is empty the inclusion is immediate.
Theory inclusion reverses set inclusion, so
\(R_{A\land B}\subseteq\operatorname{Th}(m_A\cap[B])
=R_A+B\).
For Subexpansion assume \(\neg B\notin R_A\).
Then \(m_A\) cannot be empty, since the theory of the empty
valuation set contains all formulas, and its sole member must satisfy B. That
member is also the least A-and-B valuation, giving
\(m_{A\land B}=m_A=m_A\cap[B]\).
Thus \(R_A+B=R_{A\land B}\), which proves the required
conditional inclusion. All eight cases are now checked.

Both orders tie no distinct valuations, put the only K-model
first, and put it strictly before every non-K-model; thus they
also satisfy the faithfulness definition in P8. For the input
\(\neg q\), its models are exactly \((1,0),(0,0)\).
The first order selects \((1,0)\): its theory contains p,
\(\neg q\) and \(\neg(p\to q)\). The second selects
\((0,0)\): its theory contains \(\neg p,\neg q,p\to q\).
These are opposite repairs with the same proved normative
postulate status for the fixed prior.
Conversely, modify either operation to return K whenever
\([A]=[\neg p]\), leaving all other input results unchanged.
It makes the very same observed repair on \(\neg q\), whose
valuation set differs from \([\neg p]\), but violates Success
on \(\neg p\), since \(\neg p\notin K\). Therefore a
single such choice cannot establish all eight conditions either.

At the finite-base level, retaining \(\neg q\) and either
of the original sentences is consistent, while retaining both
is inconsistent. Both repairs therefore delete one base sentence.
For the spatial contrast, strict positions A=0,B=1,C=2
satisfy A-left-B and B-left-C. Relocating A to 3 gives
B--C--A, preserves the second premise and the new A-right-C;
relocating C to -1 gives C--A--B, preserves the first premise
and the same new fact. Both move one outer object and abandon
one original ordering sentence. These stipulated equal counts
do not determine the preferred repair or prove a cognitive mechanism.
\end{proof}
\end{annexitem}
